\section{The cap achievement game on \texorpdfstring{$\mathrm{AG}(n,q)$}{AG(n,q)}}
\label{sec:affine-cap}

Let $\AG(n,q)$ denote the affine space $\Fq^n$, where $q$ is a prime power. A
\emph{cap} is a set of points containing no three distinct collinear points. The cap
achievement game starts from the empty cap and alternately adds points while preserving
the cap condition.

\begin{theorem}[Affine cap outcome law]
\label{thm:affine-cap}
For every $n\geq 1$ and every prime power $q$, the cap achievement game on $\AG(n,q)$
is a second-player win.
\end{theorem}

For $q=3$ this is the cap-set achievement game on $\mathbb F_3^n$, so the theorem settles
that game in every dimension without computation.

We use two elementary facts about affine mirrors. A \emph{cap automorphism} is a
permutation of $\Fq^n$ taking affine lines to affine lines; every invertible affine map is
one. Such a map preserves caps and legal moves.

\begin{lemma}[Line-parity lemma]
\label{lem:affine-line-parity}
Let $\sigma$ be a cap automorphism, let $A$ be a $\sigma$-symmetric cap, and let $\ell$
be a line with $\sigma(\ell)=\ell$. Suppose that on $\ell$, the map $\sigma$ acts
fixed-point-freely away from a set $F\subseteq \ell$, and that $A\cap F=\varnothing$.
If some $y\in \ell$ is a legal move at $A$, then $A\cap \ell=\varnothing$.
\end{lemma}

\begin{proof}
The set $A\cap\ell$ is $\sigma$-invariant and misses $F$, so it is partitioned into
2-element $\sigma$-orbits. Hence $|A\cap\ell|$ is even. Since $A\cup\{y\}$ is a cap and
$y\in\ell$, the line $\ell$ contains at most two points of $A\cup\{y\}$, so
$|A\cap\ell|\leq 1$. The only even possibility is $0$.
\end{proof}

\begin{lemma}[Mirror step]
\label{lem:affine-mirror-step}
Let $\sigma$ be an involutive cap automorphism, let $A$ be a $\sigma$-symmetric cap, and
let $y$ be a legal move at $A$ with $\sigma(y)\neq y$. Let $\ell$ be the line through
$y$ and $\sigma(y)$. If $A\cap\ell=\varnothing$, then $\sigma(y)$ is a legal reply and
$A\cup\{y,\sigma(y)\}$ is a cap.
\end{lemma}

\begin{proof}
If $\sigma(y)\in A$, then $y=\sigma(\sigma(y))\in \sigma(A)=A$, contradicting that $y$
is a move. Suppose $A\cup\{y,\sigma(y)\}$ contains three distinct collinear points. Any
such triple must contain $\sigma(y)$, since $A\cup\{y\}$ is already a cap. If the other
two points are in $A$, applying $\sigma$ gives a collinear triple in $A\cup\{y\}$, a
contradiction. If the triple is $\{y,\sigma(y),p\}$ with $p\in A$, then $p$ lies on
the line $\ell$, contradicting $A\cap\ell=\varnothing$.
\end{proof}

\begin{proof}[Proof of Theorem~\ref{thm:affine-cap}]
First suppose $q$ is even. If $q=2$, every subset is a cap and the board has $2^n$ points,
an even number, so the second player wins by any fixed pairing of the board.

Now let $q\geq 4$ be even. Choose a nonzero vector $v\in\Fq^n$ and define
\[
  \tau_v(x)=x+v.
\]
In characteristic $2$, $\tau_v$ is a fixed-point-free involution and an affine cap
automorphism. Let
\[
  \mathcal P=\{A\subseteq\Fq^n: A\text{ is a cap and }A=\tau_v(A)\}.
\]
The empty set lies in $\mathcal P$. If $A\in\mathcal P$ and $y$ is a legal move, let
$\ell=\{y+sv:s\in\Fq\}$ be the line through $y$ and $\tau_v(y)$. The translation
$\tau_v$ preserves $\ell$ and acts on it as $s\mapsto s+1$, fixed-point-freely. By
Lemma~\ref{lem:affine-line-parity}, $A\cap\ell=\varnothing$; by
Lemma~\ref{lem:affine-mirror-step}, $\tau_v(y)$ is a legal reply and the resulting cap
is again in $\mathcal P$. Lemma~\ref{lem:sumfree-responder} proves that the empty
position is $P$.

Now suppose $q$ is odd. The board has odd size, so a whole-board fixed-point-free mirror
cannot exist. Instead the second player builds a blocked fixed point. Let the first
player open at $a$. The second player chooses any $b\neq a$, which is legal because two
points always form a cap. Let
\[
  c=\frac{a+b}{2}.
\]
Then $a,b,c$ are distinct collinear points, so once $a$ and $b$ have been played, $c$ is
not legal for the rest of the game. Let
\[
  \sigma_c(x)=2c-x
\]
be the point reflection through $c$. This is an affine involution, fixes exactly $c$, and
swaps $a$ and $b$. Define
\[
  \mathcal P=\{A\subseteq\Fq^n: A\text{ is a cap},\ A=\sigma_c(A),\ \{a,b\}\subseteq A,
    \ c\notin A\}.
\]
The position $\{a,b\}$ lies in $\mathcal P$. If $A\in\mathcal P$ and $y$ is legal, then
$y\neq c$, since $c$ is blocked by $a,b$. Thus $\sigma_c(y)\neq y$. The line $\ell$
through $y$ and $\sigma_c(y)$ passes through $c$ and is $\sigma_c$-invariant; on
$\ell\setminus\{c\}$, the reflection is fixed-point-free. Since $A\cap\{c\}=\varnothing$,
Lemma~\ref{lem:affine-line-parity} gives $A\cap\ell=\varnothing$, and
Lemma~\ref{lem:affine-mirror-step} makes $\sigma_c(y)$ a legal reply. The new position
is again in $\mathcal P$. Therefore $\{a,b\}$ is a $P$-position by
Lemma~\ref{lem:sumfree-responder}. Since the second player can reach such a position
after every opening $a$, the empty position is $P$.
\end{proof}

The two cases explain why affine space is easier than the projective analogue. In even
characteristic, translations give a fixed-point-free pairing of the whole board. In odd
characteristic, the unavoidable fixed point of a reflection is placed on the opening secant,
where it is permanently blocked. Projective spaces have neither translations nor such a
one-point fixed set, which is why the projective cap game remains open.
